\section{Fixed-fold evaluator controls: codon frequencies versus simple length and GC}
The earlier evaluator comparisons leave a practical question open: is the codon-frequency ridge mainly exploiting simple sequence length and GC composition? A fresh frozen local plan adds those controls without changing the gene pool or selecting an optimizer result. We retain all 3,745 matched MG1655 coding-sequence rows at least 90 nucleotides long with positive PaxDb abundance. The response is log$_{10}$ ppm abundance. It is a measured abundance label from this join, not expression measured after synonymous optimization. All evaluator views use the same five shuffled outer folds (seed 7), with thirteen ridge penalties from $10^{-3}$ through $10^3$ selected by three inner training folds. No held-out target enters fitting or penalty selection.

The three fixed views are 61 sense-codon frequencies, a two-feature control comprising log$_{10}$ sequence length and GC fraction, and the concatenation of both. Codon frequencies use the same counting and frame convention as the archived nested-CV experiment. We do not tune feature families after observing the new scores, retrain a CNN, regenerate transformer embeddings or optimize any benchmark sequence. The comparison asks whether these simple controls predict the same abundance labels and whether they add predictive information under this exact estimator and data split.

We store one out-of-fold prediction per CDS row for each view, then calculate pooled Pearson correlation and ordinary RMSE across all CDS rows. Pooled correlation is not the arithmetic mean of five fold correlations. Different fold means and scales contribute to its value. The codon-only pooled value .65218 is therefore not a replacement for the historical mean-fold .65262; its five individual correlations reproduce the historical values to their recorded rounding. Preserving both prevents a metric-definition change from masquerading as a new biological finding.

\begin{table}[ht]\centering\small
\caption{Fixed nested-fold OOF predictions on 3,745 CDS rows. Intervals resample CDS rows with the fitted predictions frozen; they do not include training variability or a new external population.}
\begin{tabular}{lrrr}\hline
View & pooled Pearson $r$ & RMSE(log$_{10}$ ppm) & conditional $r$ interval\\\hline
Codon frequencies & .65218 & 1.02917 & [.63264,.67049]\\
Length and GC & .15287 & 1.34166 & [.12093,.18385]\\
Codon, length and GC & .66044 & 1.01940 & [.64116,.67843]\\\hline
\end{tabular}\end{table}

The length/GC control has much lower correlation than codon frequencies under these folds, so the codon ridge's entire score cannot be reproduced by those two features alone. Adding the controls to codon frequencies improves pooled correlation by .00826 and reduces RMSE by .00977 log$_{10}$ ppm. A paired CDS-row-index bootstrap, using 5,000 seeded resamples of the same frozen prediction ledger for all views, gives a correlation-difference interval [.00453,.01232]. This is a small conditional comparison, not proof that length or GC causes abundance changes. The GC feature is also related to the codon-frequency composition, so the combined view is not an orthogonal biological information partition.

The bootstrap sampling unit here is a CDS-row index, not a unique gene. All five trained models remain fixed within every resample. Genes share training folds and may be related through operons, sequence similarity or evolutionary history, so independence cannot be assumed as a biological fact. These intervals describe sampling of the observed prediction pairs under an iid-index approximation. They do not estimate retraining variability, train/test homology leakage, assay-specific uncertainty, independent culture replication or the effect of synthetic sequence edits. The earlier fold-bootstrap interval is also retained, with its overlapping-training-set limitation, rather than silently relabeled as a better population bound.

The fixed-fold controls establish a limited software/data result: codon frequencies predict this abundance join better than length/GC alone, and adding these two simple descriptors yields a modest incremental score. They do not validate the learned optimizer's ranking of synthesized sequences. The optimizer is evaluated outside the natural-CDS distribution after intentional selection against its own score, while these OOF predictions concern unoptimized natural genes. Shared abundance labels can make distinct evaluators agree without independent experimental confirmation. A new expression assay or an external dataset with justified transfer would be needed for a biological optimization claim.

The plan and source-join hash, all fold penalties and metrics, bootstrap settings and summary are in \path{results/oof_evaluator_audit.json}; the exact locus-tag/target/prediction ledger is \path{results/oof_evaluator_predictions.jsonl}, whose SHA-256 is retained. Run \path{scripts/oof_evaluator_audit.py}. Tests verify 3,745 rows and 3,736 unique loci, ledger/plan hashes, direct pooled-metric reproduction and the historical codon fold values. No earlier numeric result is overwritten. This experiment contributes actual model-fitting methods, measured-label results and limitations to the research body, not a new tool count or extra appendix rows.

\subsection{A failed uniqueness check exposes alternative-CDS leakage}
The first prediction-ledger regression expected 3,745 unique loci and failed: the pool contains 3,736 distinct locus tags. Eight loci have multiple distinct CDS sequences with the same abundance target: mrcB, dnaX, copA, cobB, cheA, clpB, infB and mcrB. Seven have two CDS records and infB has three, making nine rows beyond the unique-locus count. Different lengths and sequence hashes show these are alternate CDS records, not identical bytes mistakenly counted twice. Nevertheless, the shared abundance label and locus identity require grouping for a gene-level validation claim.

Six of these eight loci are split across more than one outer fold in the historical shuffled row split. A held-out CDS can therefore have an alternate CDS from the same locus in training. This is a concrete unit-of-validation failure, not a claim that it caused most of the recorded score. The diagnostic row-level results above remain unchanged, with their scope explicitly reduced. We do not simply delete the nine records or choose a representative after seeing performance.

A separate plan frozen after discovering duplicates but before the grouped outcomes shuffles the unique locus tags (seed 7) into five outer folds. All CDS records of each locus stay together. Three inner folds similarly group only the training loci, with the outer-fold index as shuffle seed. Penalty grid and model views stay fixed. Tests require zero shared locus identity across both outer and inner partitions. This controls alternative-CDS overlap, not sequence homology across different loci or shared operon identity. Those broader leakage mechanisms remain untested.

At scoring, predictions for alternate CDS rows of one held-out locus are averaged, and the single shared abundance label is used once. Thus every one of 3,736 loci receives one vote rather than a locus with more annotations receiving greater evaluation weight. Training still contains all available CDS rows within each training locus; this asymmetry is stated rather than hidden. The comparison is a new grouped protocol, not a byte-identical recreation of the historical folds. Its change cannot be attributed to leakage removal alone because the grouped partitions also differ.

\begin{table}[ht]\centering\small
\caption{Locus-grouped nested validation, scoring once per 3,736 loci after averaging alternate-CDS predictions. Fixed-prediction locus-bootstrap intervals are conditional, not independent assay or training uncertainty.}
\begin{tabular}{lrrr}\hline
View & pooled Pearson $r$ & RMSE(log$_{10}$ ppm) & conditional $r$ interval\\\hline
Codon frequencies & .65274 & 1.02876 & [.63349,.67166]\\
Length and GC & .15378 & 1.34182 & [.12257,.18566]\\
Codon, length and GC & .66170 & 1.01816 & [.64306,.68031]\\\hline
\end{tabular}\end{table}

The compatible grouped comparison retains the qualitative feature-control result: length/GC alone is weak, while the combined view yields a modest increment. Combined-minus-codon pooled correlation is .00896 with paired fixed-prediction locus-bootstrap interval [.00516,.01281]. The grouped codon value .65274 is close to the earlier row-level .65218, so this audit does not demonstrate a large score inflation from these eight loci. It demonstrates that the stated validation unit was wrong and repairs that unit in a separate experiment. A small measured effect of this repair does not justify the earlier leakage.

The 5,000 bootstrap draws now resample locus identities rather than CDS rows. Models remain frozen, and loci can still be dependent biologically. The grouped ledger records locus, number of CDS rows, its single outer-fold assignment, abundance target and averaged OOF predictions. \path{results/locus_grouped_evaluator_audit.json} retains alternate lengths, distinct-sequence counts and the exact old-fold overlap audit; \path{scripts/locus_grouped_evaluator_audit.py} regenerates the grouped results from the frozen plan. All earlier values, including the failed row-count assumption, stay visible. This is model-validation research, not evidence that synonymous designs improve measured expression or a new cancer-treatment biomarker.

\subsection{Repeated nested partitions test selection sensitivity, not a new population}
The conditional bootstrap above leaves every fitted model and selected penalty fixed. A separate frozen experiment therefore repeats the complete grouped outer and inner selection under three new outer seeds, 17, 29 and 43. Each seed has five locus-grouped outer folds and three grouped inner training folds, with inner shuffle seed equal to one hundred times the outer seed plus fold index. The same thirteen penalties and three feature views are refitted in every outer fold, giving 45 selected outer models. Scoring still gives one target and mean alternate-CDS prediction per locus.
\begin{table}[ht]\centering
\caption{Full nested refits on the same 3,736 loci. Three partition outcomes are retained; their range is not a confidence interval or independent assay replication.}
\begin{tabular}{lrrrr}\hline
Outer seed & codon $r$ & length/GC $r$ & combined $r$ & combined minus codon\\\hline
17&.65500&.15113&.66300&.00800\\
29&.65509&.15690&.66376&.00867\\
43&.65472&.15437&.66327&.00854\\\hline
\end{tabular}\end{table}
The modest combined-view increment persists in these three refitted partitions, with correlation differences .007997 through .008670. This is more than resampling already fitted predictions: training membership and penalty selection actually change. Yet all results share the same genes and abundance measurements, and three chosen partition seeds do not define a sampling distribution. We report a sensitivity range, not a training-uncertainty confidence interval or externally validated biological gain. Related loci and operons remain ungrouped beyond exact locus identity. Plans, all 45 penalties, 11,208 locus prediction records and hashes are retained in \path{results/nested_partition_audit.json} and \path{results/nested_partition_predictions.jsonl}.

\subsection{External-source scouting exposes assay and label mismatches}
The response here is PaxDb protein abundance, not mRNA expression or a measured effect of synonymous edits. A source audit identifies possible external assays but does not score them. The T7-expression experiments in the 2016 Nature article \url{https://www.nature.com/articles/nature16509} involve expression constructs and mRNA/protein analyses, requiring construct-sequence, strain, promoter and overlap checks before transfer. The 2022 nutrient-condition study \url{https://journals.plos.org/ploscompbiol/article?id=10.1371/journal.pcbi.1010641} supplies RNA-seq and ribosome profiling through GEO GSE182100. Its relative translation efficiency is not protein ppm, and repeated loci across conditions are not independent held-out genes. No external score is claimed.

For the pancreatic lane, the operator's Australian ICGC portal \url{https://pancreasexpression.org/analytics/cohort/icgc_paca_au/} labels its CDKN2A column as missense mutation. That is narrower than all mutation or genomic-alteration definitions and lacks the full burden/pathway feature panel. A truncated rendered cohort page cannot support a full denominator or patient-independent AUC. The primary pancreatic-genomics article \url{https://www.nature.com/articles/bjc2017209} discusses CDKN2A homozygous deletions, emphasizing why alteration classes must not be collapsed. Complete files, patient mapping, assay definitions and overlap checks remain open. The source ledger in \path{research/external_validation_scout.md} preserves five inspected pages, a failed official GEO fetch and these restrictions; it is preparation, not independent validation or a reason to widen the present claims.
